% Generated draft: requires mathematical and source review.
\documentclass[11pt]{article}
\usepackage[margin=25mm]{geometry}
\usepackage{fontspec}
\setmainfont{Latin Modern Roman}
\usepackage{amsmath,amssymb,mathtools,unicode-math}
\setmathfont{Latin Modern Math}
\usepackage{enumitem,needspace,xurl,hyperref}
\setlength{\parindent}{0pt}
\setlength{\parskip}{5pt}
\setlength{\emergencystretch}{3em}
\newcommand{\N}{\mathbb{N}}
\newcommand{\Z}{\mathbb{Z}}
\newcommand{\Q}{\mathbb{Q}}
\newcommand{\R}{\mathbb{R}}
\newcommand{\C}{\mathbb{C}}
\newcommand{\F}{\mathbb{F}}
\newcommand{\A}{\mathbb{A}}
\newcommand{\PP}{\mathbb{P}}
\newcommand{\E}{\mathbb{E}}
\newcommand{\eps}{\varepsilon}
\newcommand{\dd}{\,\mathrm{d}}
\begin{document}
\section*{UnsolvedMath ID: 800569}
\textbf{Source problem code:} catalogue entry
\subsection{Definitions and mathematical statement}
The supplied source asks for an asymptotically tight bound on the number of bits used by an optimal covering set of a visibility graph, together with a polynomial-time construction attaining that bound. The excerpt does not define a covering set, specify the class of visibility graphs, or specify how covering sets are encoded. These omissions prevent a complete mathematical formulation.

The optimization target can be described conditionally as follows. Once the intended instance class and representations have been specified, let \(I\) denote an instance representing a visibility graph, let \(\mathcal{C}(I)\) denote its admissible covering sets, and let \(\ell_I(C)\in\Z_{\geq 0}\) denote the number of bits in the prescribed representation of \(C\in\mathcal{C}(I)\). For instances with \(\mathcal{C}(I)\neq\varnothing\), define
\[
 b(I)=\min_{C\in\mathcal{C}(I)}\ell_I(C).
\]
The requested tasks are to determine a tight asymptotic bound for this optimum, with respect to the intended size parameter, and to construct an admissible covering set attaining that bound in polynomial time. This is a formulation schema: the supplied excerpt does not determine \(\mathcal{C}(I)\), \(\ell_I\), the size parameter, or the computational input representation. It also does not specify whether the asymptotic bound concerns each instance or the worst case over instances of a given size. No particular choices for these missing specifications are asserted here.
\subsection{Short English statement}
Problem 43 asks how many bits an optimal covering set of a visibility graph requires, up to asymptotic order, and asks for a polynomial-time construction matching that bound. The supplied excerpt leaves the covering-set definition, graph class, encoding, size parameter, and computational model unspecified; these require clarification before the problem has a precise complete statement.
\subsection{Sources}
\begin{itemize}
\item Supplied local UnsolvedMath registry record: permanent ID \(800569\), source problem code \texttt{catalogue entry}, title \emph{Visibility Graphs of Points, Segments and Polygons --- Problem 43}. Upstream revision: \texttt{d371ee917ac9647fbb06cd6fc529deec6f3dda82}.
\item Ghosh and Goswami, \emph{Unsolved Problems in Visibility Graphs of Points, Segments and Polygons} (2010), Problem 43. Supplied source link: \url{https://arxiv.org/abs/1012.5187}. The record attributes the problem to Ghosh and Goswami and gives the proposed year as 2010.
\item The supplied provenance records access on 2026-10-16 and extraction by \texttt{arxiv\_latex}. These are supplied provenance details; no independent source visit or verification is asserted. The registry labels the problem \texttt{open}; that label is supplied curation, not an independently verified assessment of its mathematical status.
\end{itemize}
\end{document}
